\section{Related Work}
\paragraph{Visual dexterous manipulation.}
ViserDex uses monocular visual pose estimation and recurrent control for dexterous reorientation, with Gaussian representations supporting visual training~\cite{bhardwaj2026viserdex}. We reuse this perception-control decomposition and adapt a released Revo3 policy to semantic goals. Closely related, GoalDex conditions a visuomotor policy on current and goal images and uses a VLM to select goal images from language~\cite{bhardwaj2026goaldex}. It studies image-goal reorientation and multi-object sim-to-real transfer. REVEAL instead addresses online semantic search when the requested marked face may be occluded: the agent obtains new evidence and binds the target during manipulation, while local control continues during the external request. GoalDex also models visual-processing delays; our distinction is this decision and control interface, not asynchronous computation itself. ManipTrans studies residual learning for dexterous bimanual transfer~\cite{li2025maniptrans}. Our residual parameterization alone is not a novelty claim.

\paragraph{Language and visual agents.}
Magma connects multimodal representations to agent actions~\cite{yang2025magma}. VERA translates visual plans into robot actions and includes language-conditioned dexterous in-hand reorientation~\cite{li2026vera}. Its hand experiments include directional rotation prompts. Consequently, language-guided in-hand rotation is not an unoccupied task category. Our setting emphasizes identifying a requested marked face under occlusion, maintaining contact during external reasoning, and checking a canonical target at autonomous stopping. Those distinctions require direct evidence, rather than a claim of being the first language-controlled hand.

\paragraph{Asynchronous robot policies.}
Real-time chunk execution~\cite{black2025rtc}, online corrections for VLA chunks~\cite{sendai2025a2c2}, and latency-aware RL~\cite{zhu2026arli} address inference delays through different execution and learning designs. DEFLECT studies delay-robust asynchronous VLAs through temporal counterfactual preferences~\cite{zhu2026deflect}. LiMA combines sparse semantic intention and asynchronous action generation for dexterous manipulation~\cite{chen2026lima}. These are close methodological references. Our current implementation learns small teacher-supervised controllers around a separate recurrent policy and public semantic tools. A controlled comparison with a generic time-aware residual policy is still needed to show that the active semantic goal and request anchor add value beyond reactive correction.
